\subsection{Proof of Theorem 1: second part}

Consider again the exact sequence

\[
0 \longrightarrow t ^ {r} / t ^ {r + 1} \xrightarrow {i _ {r}} A _ {0} / t ^ {r + 1} \xrightarrow {p _ {r}} A _ {0} / t ^ {r} \longrightarrow 0\tag{16}
\]

and the associated exact sequence of torsion modules

\[
\begin{array}{l} \text {Tor} _ {1} ^ {\mathrm{A} _ {0}} (\mathrm{A} _ {0} / t ^ {r + 1}, \mathrm{M}) \longrightarrow \mathrm{Tor} _ {1} ^ {\mathrm{A} _ {0}} (\mathrm{A} _ {0} / t ^ {r}, \mathrm{M}) \\ \longrightarrow (t ^ {r} / t ^ {r + 1}) \otimes_ {\mathrm{A} _ {0}} \mathrm{M} \xrightarrow {i , \otimes 1 _ {\mathrm{M}}} (\mathrm{A} _ {0} / t ^ {r + 1}) \otimes_ {\mathrm{A} _ {0}} \mathrm{M} \xrightarrow {p , \otimes 1 _ {\mathrm{M}}} (\mathrm{A} _ {0} / t ^ {r}) \otimes_ {\mathrm{A} _ {0}} \mathrm{M} \longrightarrow 0. \end{array} \tag {18}
\]

The kernel of \(p_r \otimes 1_M\) is identified with \(x^r M/x^{r+1} M\) ; moreover, \(t^r/t^{r+1}\) is annihilated by the \(T_i\) and is identified with the homogeneous component of degree \(r\) of \(A_0\) , so that the homomorphism \((t^r/t^{r+1}) \otimes_{A_0} M \to x^r M/x^{r+1} M\) deduced from \(i_r \otimes 1_M\) is identified with the homogeneous component of degree \(r\) of the homomorphism \(\beta_M^x\) . It therefore follows from the exact sequence (18) that condition (v) is equivalent to

(v'): the homomorphism \(\operatorname{Tor}_{1}^{\mathbf{A}_{0}}(p_{r},1_{\mathbf{M}})\) : \(\operatorname{Tor}_{1}^{\mathbf{A}_{0}}(\mathbf{A}_{0}/t^{r+1},\mathbf{M})\to\operatorname{Tor}_{1}^{\mathbf{A}_{0}}(\mathbf{A}_{0}/t^{r},\mathbf{M})\) is surjective for all \(r\geqslant1\) .

It remains for us to prove the implication (v) \(\Rightarrow\) (i) when the modules

\[
\mathbf {M} / (x _ {1} \mathbf {M} + \dots + x _ {i - 1} \mathbf {M})
\]

are separated for the \(\mathfrak{x}\) -adic topology ( \(1 \leqslant i \leqslant n\) ). Let \(\overline{\mathbf{M}}\) be the A-module \(\mathbf{M}/x_1 \mathbf{M}\) . By definition, the sequence \(\mathbf{x}\) is regular for \(\mathbf{M}\) if and only if \((x_1)_\mathbf{M}\) is injective and the sequence \(\mathbf{x}' = (x_2, \ldots, x_n)\) is regular for \(\overline{\mathbf{M}}\) . Reasoning by induction on \(n\) , it therefore suffices to prove that, if \(\mathbf{M}\) is separated for the \(\mathfrak{x}\) -adic topology and if \(\beta_{\mathbf{M}}^{\mathbf{x}}\) is bijective, then \((x_1)_\mathbf{M}\) is injective and \(\beta_{\mathbf{M}}^{\mathbf{x}'}\) is bijective. Now the bijectivity of \(\beta_{\mathbf{M}}^{\mathbf{x}}\) implies in particular that the homothety with ratio \(x_1\) to \(\bigoplus_{r} x^r \mathbf{M}/x^{r+1} \mathbf{M}\) is injective, hence that \(\text{Ker } (x_1)_\mathbf{M} \subset \bigcap_{i} x^i \mathbf{M}\) and consequently that \((x_1)_\mathbf{M}\) is injective if the \(\mathfrak{x}\) -adic topology on \(\mathbf{M}\) is separated.

We are thus reduced to proving that if \((x_{1})_{\mathbf{M}}\) is injective and if M satisfies the condition \((\mathbf{v}')\) , then \(\overline{M}\) satisfies the condition \((\mathbf{v}')\) relative to the sequence \(x'\) .

By hypothesis, we have an exact sequence

\[
0 \longrightarrow \mathrm{M} \xrightarrow {(x _ {1}) _ {M}} \mathrm{M} \longrightarrow \overline {{\mathrm{M}}} \longrightarrow 0;
\]

let us set \(\overline{\mathbf{A}}_{0} = \mathbf{A}_{0}/\mathbf{T}_{1}\mathbf{A}_{0},\quad\overline{\mathbf{t}} = \mathbf{t}\overline{\mathbf{A}}_{0}\) . Let \(q:L\to M\) be a free resolution of the \(\mathbf{A}_{0}\) -module M; since the homothety of ratio \(\mathbf{T}_{1}\) is injective in \(\mathbf{A}_{0}\) , it is injective in L, and we have an exact sequence of complexes

\[
0 \longrightarrow \mathrm{L} \xrightarrow {(x _ {1}) _ {\mathrm{L}}} \mathrm{L} \longrightarrow \overline {{\mathrm{L}}} \longrightarrow 0
\]

with \(\overline{\mathbf{L}} = \mathbf{L} / x_1\mathbf{L}\) , and a commutative diagram

\[
\begin{array}{c} 0 \longrightarrow \mathrm{L} \stackrel {(x _ {1}) _ {\mathrm{L}}} {\longrightarrow} \mathrm{L} \longrightarrow \overline {{\mathrm{L}}} \longrightarrow 0 \\ \Big \downarrow^ {q} \quad \Big \downarrow^ {q} \quad \Big \downarrow^ {\overline {{q}}} \\ 0 \longrightarrow \mathrm{M} \stackrel {(x _ {1}) _ {\mathrm{M}}} {\longrightarrow} \mathrm{M} \longrightarrow \overline {{\mathrm{M}}} \longrightarrow 0. \end{array}
\]

Since \(q\) is a homologism, \(\overline{q}:\overline{\mathbf{L}}\to\overline{\mathbf{M}}\) is a free resolution of the \(\overline{\mathbf{A}}_{0}\) -module \(\overline{\mathbf{M}}\) (cor. 1, p.~30). For every \(\overline{\mathbf{A}}_{0}\) -module P, there is a canonical isomorphism

\[
\mathbf {P} \otimes_ {\mathbf {A} _ {0}} \mathbf {L} \rightarrow \mathbf {P} \otimes_ {\overline {{{\mathbf {A}}}} _ {0}} \overline {{{\mathbf {L}}}},
\]

whence, by passing to homology, an isomorphism

\[
\varphi_ {P}: \operatorname{Tor} _ {1} ^ {\mathbf {A} _ {0}} (\mathbf {P}, \mathbf {M}) \rightarrow \operatorname{Tor} _ {1} ^ {\overline {{\mathbf {A}}} _ {0}} (\mathbf {P}, \overline {{\mathbf {M}}}).
\]

If \(u: \mathbf{P} \to \mathbf{P}'\) is a \(\overline{\mathbf{A}}\) -homomorphism, then

\[
\varphi_ {\mathrm{P} ^ {\prime}} \circ \operatorname{Tor} _ {1} ^ {\mathbf {A} _ {0}} (u, 1 _ {\mathrm{M}}) = \operatorname{Tor} _ {1} ^ {\overline {{\mathbf {A}}} _ {0}} (u, 1 _ {\mathrm{M}}) \circ \varphi_ {\mathrm{P}}.
\]

That being so, suppose condition (v') is satisfied for M, and let us prove that it holds for \(\overline{M}\) . Let \(r\) be an integer \(\geqslant 1\) ; we have a commutative diagram with exact rows

\[
\begin{array}{c} 0 \longrightarrow \mathrm{A} _ {0} / t ^ {r} \xrightarrow {\mathrm{T} _ {1}} \mathrm{A} _ {0} / t ^ {r + 1} \longrightarrow \overline {{\mathrm{A}}} _ {0} / \overline {{t}} ^ {r + 1} \longrightarrow 0 \\ \downarrow^ {p _ {r - 1}} \quad \downarrow^ {p _ {r}} \quad \downarrow^ {\bar {p} _ {r}} \\ 0 \longrightarrow \mathrm{A} _ {0} / t ^ {r - 1} \xrightarrow {\mathrm{T} _ {1}} \mathrm{A} _ {0} / t ^ {r} \longrightarrow \overline {{\mathrm{A}}} _ {0} / \overline {{t}} ^ {r} \longrightarrow 0 \end{array}
\]

from which we deduce a commutative diagram with exact rows

\[
\begin{array}{c c c c c} \operatorname{Tor} _ {1} ^ {\mathbf {A} _ {0}} (\mathbf {A} _ {0} / t ^ {r + 1}, \mathbf {M}) & \longrightarrow & \operatorname{Tor} _ {1} ^ {\mathbf {A} _ {0}} (\overline {{\mathbf {A}}} _ {0} / \overline {{t}} ^ {r + 1}, \mathbf {M}) & \longrightarrow & \mathbf {M} / x ^ {r}   \mathbf {M} \xrightarrow {x _ {1}} \mathbf {M} / x ^ {r + 1}   \mathbf {M} \\ \operatorname{Tor} _ {1} (p _ {r}, 1) & \downarrow & \operatorname{Tor} _ {1} (\bar {p} _ {r}, 1) & \downarrow & \downarrow \\ \operatorname{Tor} _ {1} ^ {\mathbf {A} _ {0}} (\mathbf {A} _ {0} / t ^ {r}, \mathbf {M}) & \longrightarrow & \operatorname{Tor} _ {1} ^ {\mathbf {A} _ {0}} (\overline {{\mathbf {A}}} _ {0} / \overline {{t}} ^ {r}, \mathbf {M}) & \longrightarrow & \mathbf {M} / x ^ {r - 1}   \mathbf {M} \xrightarrow {x _ {1}} \mathbf {M} / x ^ {r}   \mathbf {M}. \end{array}
\]

Now multiplication by \(x_{1}\) defines an injection of \(\mathbf{M}/\mathfrak{x}^{r}\mathbf{M}\) to \(\mathbf{M}/\mathfrak{x}^{r+1}\mathbf{M}\) : this follows immediately, by induction on \(r\) , from the exact sequence

\[
0 \rightarrow x ^ {r - 1} M / x ^ {r} M \rightarrow M / x ^ {r} M \rightarrow M / x ^ {r - 1} M \rightarrow 0
\]

and from the injectivity of the homothety with ratio \(x_{1}\) to \(\bigoplus_{r\geqslant 0}(\mathfrak{x}^{r}\mathrm{M}/\mathfrak{x}^{r+1}\mathrm{M})\) . Condition (v) therefore implies that the homomorphism \(\operatorname{Tor}_{1}^{\mathbf{A}_{0}}(\overline{p}_{r},1_{\mathbf{M}})\) is surjective for all \(r\geqslant 1\) . By composition with the isomorphisms \((\varphi_{\overline{\mathbf{A}}_{0}/\overline{\mathbf{t}}^{r+1}})^{-1}\) and \(\varphi_{\overline{\mathbf{A}}_{0}/\overline{\mathbf{t}}^{r}}\) , we deduce that the homomorphism

\[
\operatorname{Tor} _ {1} ^ {\overline {{\mathbf {A}}} _ {0}} \left(\bar {p} _ {r}, 1 _ {\mathbf {M}}\right): \operatorname{Tor} _ {1} ^ {\overline {{\mathbf {A}}} _ {0}} \left(\overline {{\mathbf {A}}} _ {0} / \bar {t} ^ {r + 1}, \mathbf {M}\right)\rightarrow \operatorname{Tor} _ {1} ^ {\overline {{\mathbf {A}}} _ {0}} \left(\overline {{\mathbf {A}}} _ {0} / \bar {t} ^ {r}, \mathbf {M}\right)
\]

is surjective for \(r \geqslant 1\) , whence condition (v') for \(\overline{\mathbf{M}}\) . This completes the proof of the theorem.

